\section{Introducción}

\subsection{Planteamiento del problema}
El Braille organiza el texto en celdas de dos columnas por tres filas. En visión
por computadora, el reconocimiento óptico de Braille (OBR) pasó de la
umbralización y la morfología matemática \parencite{antonacopoulos2004} al
aprendizaje profundo supervisado \parencite{li2020}. Los detectores recientes
localizan la celda completa y la clasifican en 63 categorías, lo que confiere
invarianza ante rotaciones y deformaciones del papel
\parencite{ovodov2021,sun2025}. Sin embargo, se formularon para documentos
editoriales impresos o dependen de servicios en la nube; el Braille escolar
manuscrito con regleta y punzón ha recibido escasa atención.

En el Perú, el Informe Defensorial N.º 183 reportó que el Programa Presupuestal
106 representaba apenas el 0.62\,\% del presupuesto educativo y que más del
88\,\% de las personas con discapacidad en edad escolar no estaba matriculada
\parencite{defensoria2019}. En Moquegua, el docente de aula no puede leer el
cuaderno en relieve y la corrección queda retenida hasta la visita del
especialista; la retroalimentación formativa, cuyo valor depende de la
inmediatez \parencite{hattie2007}, queda anulada.

\subsection{Antecedentes y vacío}
La literatura se ordena en métodos clásicos, segmentación semántica y detección
de objetos de una etapa. \textcite{li2020} propusieron BraUNet; \textcite{ovodov2021}
plantearon la detección sobre la celda completa; \textcite{sun2025} exploraron
YOLOv11 en escenas naturales. Para el español, \textcite{fuentes2026} presentaron
una aplicación móvil con CNN, pero no evalúan Braille manuscrito escolar ni el
compromiso precisión--eficiencia en gama baja. Persisten tres vacíos: el Braille
escolar manuscrito, la ejecución sin conexión y una evaluación integrada de
precisión, concordancia, eficiencia y efecto pedagógico.

\subsection{Objetivos e hipótesis}
El objetivo general es evaluar configuraciones de una CNN ligera para OBR en
móvil e identificar la de mejor compromiso entre precisión, eficiencia y
oportunidad de la retroalimentación docente. Se plantean tres hipótesis:
YOLOv11 supera a YOLOv8n en mAP@0.5 y F1 macro (H1); INT8 reduce latencia y
memoria con una degradación acotada de mAP (H2); y el sistema reduce el tiempo de
revisión por página (H3).
